\documentclass[11pt]{amsart}
\usepackage[margin=1.15in]{geometry}
\usepackage{amsmath,amssymb,amsthm,mathtools}
\usepackage{enumitem}
\usepackage[colorlinks=true,linkcolor=blue,citecolor=blue]{hyperref}

\newtheorem{theorem}{Theorem}[section]
\newtheorem{lemma}[theorem]{Lemma}
\newtheorem{proposition}[theorem]{Proposition}
\newtheorem{corollary}[theorem]{Corollary}
\newtheorem{thmx}{Theorem}
\renewcommand{\thethmx}{\Alph{thmx}}
\theoremstyle{definition}
\newtheorem{definition}[theorem]{Definition}
\newtheorem{example}[theorem]{Example}
\newtheorem{question}[theorem]{Question}
\theoremstyle{remark}
\newtheorem{remark}[theorem]{Remark}

\newcommand{\GG}{\mathbb{G}}
\newcommand{\ZZ}{\mathbb{Z}}
\newcommand{\QQ}{\mathbb{Q}}
\newcommand{\FF}{\mathbb{F}}
\newcommand{\acl}{\operatorname{acl}}
\newcommand{\trdeg}{\operatorname{trdeg}}
\newcommand{\rk}{\operatorname{rk}}
\newcommand{\Frob}{\operatorname{Frob}}
\newcommand{\lev}{\operatorname{lev}}
\newcommand{\hgt}{\operatorname{ht}}
\newcommand{\lt}{\operatorname{lt}}
\newcommand{\chr}{\operatorname{char}}
\newcommand{\dlog}{\operatorname{dlog}}
\newcommand{\sep}{\mathrm{sep}}
\newcommand{\sm}{\smallsetminus}
\newcommand{\la}{\langle}
\newcommand{\ra}{\rangle}
\newcommand{\KM}{K^{M}}
\newcommand{\Kr}{\mathcal{K}}

\title[Milnor $K$-theory and field embeddings]{Steinberg-compatible homomorphisms of multiplicative groups\\ and the reconstruction of fields from Milnor $K$-theory}
\date{November 18, 2026. Working draft.}

\begin{document}

\begin{abstract}
Let $K$ and $L$ be fields and let $\psi$ be a homomorphism from $K^\times$ to the rational multiplicative group of $L$ modulo constants.
We show that if $\psi$ is injective modulo constants, if $\psi(x)$ and $\psi(1-x)$ are always algebraically dependent, and if the image of $\psi$ contains
two (in positive characteristic: three) algebraically independent elements, then $\psi$ is a rational multiple of the map induced by an embedding of $K$ into the
perfect closure of $L$. The proof is elementary: it uses K\"ahler differentials, separability and, in positive characteristic, a book-keeping of Frobenius twists.
As a consequence, every homomorphism between rational Milnor $K$-rings of fields which is injective in degree one and non-degenerate comes from an embedding of fields;
in particular a field of absolute transcendence degree at least $3$ (at least $2$ in characteristic $0$) is determined up to purely inseparable extensions by $\QQ\otimes\KM_*$.
\end{abstract}

\maketitle
\setcounter{tocdepth}{1}
\tableofcontents

\section{Introduction}

\subsection{Statement for multiplicative groups}
Let $L$ be a field and $l\subset L$ a relatively algebraically closed subfield. Put
\[ V_{L|l}=(L^\times/l^\times)\otimes_\ZZ\QQ , \]
a $\QQ$-vector space which we write additively; the class of $a\in L^\times$ is denoted $[a]$.
Every non-zero $v\in V_{L|l}$ satisfies $Nv=[a]$ for some integer $N\ge1$ and some $a\in L\sm l$, and elements $v_1,\dots,v_m$ are called \emph{algebraically independent over $l$}
if the corresponding $a_1,\dots,a_m$ are; this does not depend on choices. The group does not change if $L$ and $l$ are replaced by their perfect closures $L^i$ and $l^i$.

\begin{thmx}\label{thm:A}
Let $K$ and $L$ be fields, $k\subset K$ a subfield, $l\subset L$ a relatively algebraically closed subfield, and let $\psi\colon K^\times\to V_{L|l}$ be a homomorphism such that
\begin{enumerate}[label=\textup{(A\arabic*)}]
\item\label{A1} $\ker\psi=k^\times$;
\item\label{A2} for every $x\in K\sm k$, the elements $\psi(x)$ and $\psi(1-x)$ are algebraically dependent over $l$;
\item\label{A3} there are $x_1,\dots,x_d\in K$ with $\psi(x_1),\dots,\psi(x_d)$ algebraically independent over $l$, where $d=2$ if $\chr L=0$ and $d=3$ if $\chr L>0$;
\item\label{A4} if $\chr L>0$, then $\chr K>0$.
\end{enumerate}
Then $\chr K=\chr L$, the field $k$ is relatively algebraically closed in $K$, and there are a field embedding $\tau\colon K\to L^i$ with $\tau^{-1}(l^i)=k$ and a number $r\in\QQ^\times$ such that
\[ \psi(x)=r\cdot[\tau(x)]\qquad\text{for all }x\in K^\times . \]
If $(\tau',r')$ is another such pair, then $\tau'=\Frob^j\circ\tau$ and $r'=p^{-j}r$ for some integer $j$, where $p$ is the characteristic exponent; in characteristic $0$ the pair is unique.
\end{thmx}

Conversely, every pair $(\tau,r)$ as in the theorem defines a homomorphism $\psi=r[\tau]$ satisfying \ref{A1} and \ref{A2}.
For function fields of transcendence degree $\ge2$ over algebraically closed fields, a statement of this kind is due to Bogomolov and Tschinkel \cite{BT}.
Theorem~\ref{thm:A} has no assumption on the fields, and its proof is different.

\subsection{Statements for Milnor \texorpdfstring{$K$}{K}-theory}
For a field $K$ let $\KM_*(K)$ be its Milnor ring, the quotient of the tensor algebra of $K^\times$ by the ideal generated by the elements $x\otimes(1-x)$, and put $\Kr_*(K)=\QQ\otimes\KM_*(K)$.
For a subfield $k\subset K$ let $\Kr_*(K|k)$ be the quotient of $\Kr_*(K)$ by the ideal generated by $k^\times\subset\Kr_1(K)$; if $k$ is relatively algebraically closed, $\Kr_1(K|k)=V_{K|k}$.
A field embedding $\tau\colon K\to L$ with $\tau(k)\subset l$ induces $\tau_*\colon\Kr_*(K|k)\to\Kr_*(L|l)$, and for $r\in\QQ^\times$ we write $r^{\deg}$ for the automorphism of a graded $\QQ$-algebra which is multiplication by $r^n$ in degree $n$.
To keep the statements simple we assume in Theorems~\ref{thm:B}--\ref{thm:D} that the fields are perfect; this is harmless, because $\Kr_*$ does not change under purely inseparable extensions.

\begin{thmx}[Relative version]\label{thm:B}
Let $K,L$ be perfect fields with relatively algebraically closed subfields $k,l$, and let $\varphi\colon\Kr_*(K|k)\to\Kr_*(L|l)$ be a homomorphism of graded $\QQ$-algebras.
Assume that $\varphi_1$ is injective, that $\psi=\varphi_1\circ[\,\cdot\,]$ satisfies \ref{A3}, and that \ref{A4} holds.
Then there are a field embedding $\tau\colon K\to L$ with $\tau^{-1}(l)=k$ and $r\in\QQ^\times$ such that $\varphi=r^{\deg}\circ\tau_*$. The pair $(\tau,r)$ is unique up to $(\Frob^j\circ\tau,p^{-j}r)$.
\end{thmx}

\begin{thmx}[Absolute version]\label{thm:C}
Let $K,L$ be perfect fields, and let $k_0,l_0$ be the algebraic closures of the prime fields in $K$ and $L$. Assume that $K\ne k_0$.
Let $\varphi\colon\Kr_*(K)\to\Kr_*(L)$ be a homomorphism of graded $\QQ$-algebras such that $\varphi_1$ is injective, the composition $K^\times\to\Kr_1(L)\to V_{L|l_0}$ satisfies \ref{A3}, and \ref{A4} holds.
Then there are a field embedding $\tau\colon K\to L$ and $r\in\QQ^\times$ such that $\varphi=r^{\deg}\circ\tau_*$, and $\tau^{-1}(l_0)=k_0$. The pair $(\tau,r)$ is unique up to $(\Frob^j\circ\tau,p^{-j}r)$.
\end{thmx}

\begin{thmx}[Isomorphisms]\label{thm:D}
Let $K$ and $L$ be fields each of which contains three elements that are algebraically independent over the prime field; if one of them has characteristic $0$, two such elements suffice.
Then every isomorphism of graded $\QQ$-algebras $\varphi\colon\Kr_*(K)\to\Kr_*(L)$ is of the form $r^{\deg}\circ\tau_*$ for an isomorphism of fields $\tau\colon K^i\to L^i$ and some $r\in\QQ^\times$.
In particular $K^i\cong L^i$.
\end{thmx}

For transcendence degree $\ge5$, Theorem~\ref{thm:D} and the analogue of Theorem~\ref{thm:B} for isomorphisms were proved by Topaz \cite{T}; the passage from $K$-theory to fields there goes through
the classification of ``alternating pairs'' by valuations and through the reconstruction of a field from its lattice of relatively algebraically closed subfields.
The present approach avoids both: the only input from $K$-theory is the Steinberg relation $\{x,1-x\}=0$ and one family of homomorphisms out of $K_2$ (Lemma~\ref{lem:tame}).

\subsection{The idea}
In characteristic $0$ the proof of Theorem~\ref{thm:A} takes two pages.
The logarithmic derivative is an injective $\QQ$-linear map $\dlog\colon V_{L|l}\to\Omega_{L|l}$, and two elements of $V_{L|l}$ are algebraically dependent if and only if their images are proportional.
So $\Delta=\dlog\circ\psi$ is a homomorphism $K^\times\to\Omega_{L|l}$ with $\Delta(1-x)=\mu(x)\Delta(x)$ for a function $\mu$ with values in $L^\times$.
The identity $1-xy=(1-x)(1-u)$ with $u=x(1-y)/(x-1)$, applied to two elements with linearly independent $\Delta(x),\Delta(y)$, gives a functional equation for $\mu$ which says precisely that
\[ t(x)=\frac{\mu(x)}{\mu(x)-1} \]
is multiplicative; and $t(1-x)=1-t(x)$ holds for formal reasons. Hence $t$ is a field embedding $\tau$, and a short argument identifies $\psi$ with a rational multiple of $[\tau]$.

In characteristic $p$ the group $V_{L|l}$ is $p$-divisible and has no logarithmic derivative. One has to choose a non-perfect \emph{model} $L_A\subset L$ and use $\Omega_{L_A|l}$, which only sees the elements of the
correct Frobenius \emph{level}. The whole difficulty is to show that $\psi(x)$ and $\psi(1-x)$ have the same level (Proposition~\ref{prop:levels}).
This is where three independent elements and the hypothesis \ref{A4} are used: a configuration of seven points shows that the level discrepancy is a homomorphism $K^\times\to\ZZ$, which must vanish because $K^\times$ is divisible by the characteristic.
Once this is known, the characteristic $0$ computation runs in the model.

The exposition is uniform in the characteristic. A reader interested only in characteristic $0$ can skip Sections~\ref{sec:levels}--\ref{sec:quadrangle} and Section~\ref{sec:stepA}, read $L_A=L$ everywhere, and ignore all statements about levels, heights and Frobenius.

\subsection{What is used}
Besides linear algebra, the proof uses the following standard material, collected in Section~\ref{sec:conv} and Lemma~\ref{lem:tame}:
K\"ahler differentials of separably generated field extensions; Gauss' lemma and unique factorization in polynomial rings; the norm of a finite separable extension;
and, for the $K$-theoretic statements only, the existence of an extension of a valuation to a larger field. No result about Milnor $K$-theory beyond its definition is used.

\subsection{Status}
This is a working draft which has not been refereed or independently checked. The statements go beyond the literature cited (in the transcendence degree, and in allowing non-surjective homomorphisms), so the text should be read as a detailed proposal for a proof.
The characteristic $0$ part (Lemmas~\ref{lem:onevar}--\ref{lem:theta} read with $L_A=L$, and Section~\ref{sec:K}) is short and can be checked line by line. In positive characteristic the steps carrying most of the weight, and the first places to look for a mistake, are Lemmas~\ref{lem:trichotomy}, \ref{lem:quadrangle}, \ref{lem:gauge}, \ref{lem:cocycle} and Proposition~\ref{prop:rational}.

\section{Conventions and standard facts}\label{sec:conv}

\subsection{Points and flats}
Let $L$ be a field and $l\subset L$ a relatively algebraically closed subfield. A \emph{flat} of $L|l$ is a relatively algebraically closed subfield of $L$ containing $l$; its \emph{rank} is its transcendence degree over $l$.
A \emph{point} is a flat of rank $1$ and a \emph{line} is a flat of rank $2$. For $a\in L\sm l$ we write $[a]_L$, or simply $Q_a$, for the point containing $a$, the relative algebraic closure of $l(a)$ in $L$.
For points $P_1,\dots,P_m$ we write $\la P_1,\dots,P_m\ra$ for the flat they generate; they are \emph{collinear} if this flat has rank $\le2$. We write $\rk\{P_1,\dots,P_m\}$ for the rank of $\la P_1,\dots,P_m\ra$.
Two distinct points intersect in $l$, and two distinct lines through a point $P$ intersect in $P$.

\subsection{Frobenius and the group \texorpdfstring{$\Gamma$}{Gamma}}
For a perfect field $L$ let $p=p(L)$ be its characteristic exponent ($p=1$ in characteristic $0$) and $\Frob\colon x\mapsto x^p$ its Frobenius automorphism. We put $\Gamma=\Gamma(L)=\ZZ$ if $p>1$ and $\Gamma=\{0\}$ if $p=1$,
and $\Gamma_{\ge0}=\{j\in\Gamma:j\ge0\}$. Thus $j\mapsto\Frob^j$ is injective on $\Gamma$ (our fields are infinite), and all integers measuring Frobenius twists are elements of $\Gamma$.
We write $\FF$ for the prime field of the field under consideration.

\subsection{The group \texorpdfstring{$V_{L|l}$}{V}}\label{ss:V}
Let $L$ be perfect and $l\subset L$ relatively algebraically closed. We abbreviate $V=V_{L|l}$.
\begin{enumerate}[label=\textup{(V\arabic*)}]
\item\label{V1} $[a]=0$ in $V$ if and only if $a\in l^\times$, because $L^\times/l^\times$ is torsion-free.
\item\label{V2} Every $v\in V\sm0$ has a \emph{representative}: a pair $(N,a)$ with $N\ge1$ an integer prime to $p$ and $a\in L\sm l$ such that $Nv=[a]$.
Indeed $v=\frac1n[b]$ with $n=p^jn'$, and one takes $N=n'$, $a=\Frob^{-j}(b)$. If $(N',a')$ is another representative then $a^{N'}=c\,a'^{N}$ with $c\in l^\times$.
\item\label{V3} Consequently the point $Q_v=Q_a$ depends only on $v$. For a point $Q$, the set $G_Q=\{v: v=0\text{ or }Q_v=Q\}$ is a subgroup of $V$, the image of $(Q^\times/l^\times)\otimes\QQ$.
Elements $v_1,\dots,v_m\in V\sm0$ are \emph{algebraically independent} if their representatives are; $v,v'$ are dependent if and only if $Q_v=Q_{v'}$.
\end{enumerate}
For a homomorphism $\psi\colon K^\times\to V$ and $x\notin\ker\psi$ we write $Q_x=Q_{\psi(x)}$, and we call $x_1,\dots,x_m$ \emph{$\psi$-independent} if $\psi(x_1),\dots,\psi(x_m)$ are algebraically independent.


\subsection{Standard facts}
\begin{lemma}[Standard facts]\label{lem:standard}
Let $l$ be a perfect field of characteristic exponent $p$.
\begin{enumerate}[label=\textup{(S\arabic*)}]
\item\label{S1} Let $f\in l[Y_1,\dots,Y_m]$. If $\partial f/\partial Y_1=0$, then $Y_1$ does not occur in $f$, or $p>1$ and $f\in l[Y_1^p,Y_2,\dots,Y_m]$.
If all partial derivatives of $f$ vanish, then $f$ is constant, or $p>1$ and $f$ is a $p$-th power.
\item\label{S2} Let $E\supset l$ be a field which is separable algebraic over $l(A)$ for an algebraically independent set $A$. Then $\{da:a\in A\}$ is a basis of $\Omega_{E|l}$, and $d(\Frob y)=0$ for $y\in E$ if $p>1$.
If $y\in E$ is separable algebraic over $l(x_1,\dots,x_r)\subset E$, then $dy\in\sum_iE\,dx_i$; more precisely $g_Y\,dy+\sum_ig_{X_i}\,dx_i=0$ for every polynomial relation $g(x_1,\dots,x_r,y)=0$ over $l$.
\item\label{S3} In the situation of \ref{S2} assume moreover that $l$ is relatively algebraically closed in $E$, and let $y\in E\sm l$. Then there is a largest $j\in\Gamma_{\ge0}$ with $\Frob^{-j}(y)\in E$, and $d(\Frob^{-j}y)\neq0$.
In particular, if $dy=0$ then $p>1$ and $y\in\Frob(E)$.
\end{enumerate}
\end{lemma}

\begin{proof}
\ref{S1} and \ref{S2} are standard. For \ref{S3}: if $p=1$ the claim is that $dy\ne0$ for $y$ transcendental over $l$, which holds because $E$ is separable over $l(y)$.
Let $p>1$. Then the kernel of $d$ is $E^p$, so it suffices to show that $y\notin\bigcap_jE^{p^j}$.
Suppose $z_j=y^{p^{-j}}\in E$ for all $j$. Let $A_0\subset A$ be finite with $y$ algebraic over $l(A_0)$ and put $E'=l(A_0)(y)$.
Each $z_j$ is algebraic over $l(A_0)$ and separable over $l(A)$; since $l(A)$ is purely transcendental over $l(A_0)$, the minimal polynomial of $z_j$ over $l(A_0)$ stays irreducible over $l(A)$, so
$z_j$ is separable over $l(A_0)$, hence over $E'$. It is also purely inseparable over $E'$, so $z_j\in E'$.
Now $E'$ is finitely generated over $l$ and $y$ is transcendental over $l$, so there is a discrete valuation $v$ of $E'$, trivial on $l$, with $v(y)\ne0$. Then $v(z_j)=v(y)/p^j\in\ZZ$ for all $j$, a contradiction.
\end{proof}

\begin{lemma}[Eigenvalue lemma]\label{lem:R}
Let $E\supset l$ be fields, $a\in E$ transcendental over $l$, and $\delta$ a derivation of $E$ with $\delta(l)=0$ and $\delta(a)=a$.
Let $b\in E^\times$ be separable algebraic over $l(a)$ with $\delta(b)=\kappa\,b$ for some $\kappa\in l$. Then $\kappa$ lies in the prime field $\FF$ of $l$.
\end{lemma}

\begin{proof}
Let $P(Y)=\sum_{j=0}^nP_jY^j\in l[a][Y]$ be irreducible with $P(b)=0$, so the $P_j\in l[a]$ have no common factor. Then $P_0\ne0$ since $b\neq0$, and since $P$ is separable there is $j_1\geq1$ with $j_1\ne0$ in $l$ and $P_{j_1}\neq0$.
On $l[a]$ the derivation $\delta$ is $Q\mapsto aQ'$. Applying $\delta$ to $P(b)=0$ gives $R(b)=0$ for
\[ R(Y)=\sum_j\bigl(aP_j'+j\kappa P_j\bigr)Y^j . \]
As $R$ has degree $\le n$ in $Y$ and $P$ is a minimal polynomial of $b$ over $l(a)$, we get $R=\lambda P$ with $\lambda\in l(a)$. Since the $P_j$ are coprime, $\lambda\in l[a]$, and comparing degrees in $a$ shows $\lambda\in l$.
So $aP_j'=(\lambda-j\kappa)P_j$ for all $j$. If a non-zero polynomial $Q=\sum c_ia^i$ satisfies $aQ'=\varepsilon Q$ with $\varepsilon\in l$, then $(i-\varepsilon)c_i=0$ for all $i$, so $\varepsilon\in\FF$.
Applying this to $P_0$ and to $P_{j_1}$ gives $\lambda\in\FF$ and $\lambda-j_1\kappa\in\FF$, hence $\kappa\in\FF$.
\end{proof}

\section{Levels and normalized triples}\label{sec:levels}

In Sections~\ref{sec:levels}--\ref{sec:quadrangle}, $L$ is a perfect field of characteristic exponent $p$, $l\subset L$ is a relatively algebraically closed
(hence perfect) subfield, and $\Gamma=\Gamma(L)$. ``Point'' and ``line'' refer to $L|l$ as in Section~\ref{sec:conv}. Every point is stable under $\Frob^{\pm1}$.

\begin{lemma}\label{lem:sepcrit}
Let $m\ge2$ and let $y_1,\dots,y_m\in L$ be algebraically dependent over $l$, any $m-1$ of them being algebraically independent. Then:
\begin{enumerate}[label=\textup{(\alph*)}]
\item The ideal of polynomial relations among $(y_1,\dots,y_m)$ over $l$ is generated by an irreducible $f\in l[Y_1,\dots,Y_m]$ in which every variable occurs.
\item At least one partial derivative of $f$ is non-zero.
\item For each $i$ the following are equivalent: $y_i$ is separable over $l(y_j : j\ne i)$; $\partial f/\partial Y_i\neq0$; $(\partial f/\partial Y_i)(y_1,\dots,y_m)\ne0$.
\end{enumerate}
\end{lemma}

\begin{proof}
(a) The ideal is a prime of height one in a factorial ring. If $Y_i$ did not occur in $f$, the other $y_j$ would be dependent.
(b) Otherwise $f$ would be constant or a $p$-th power by \ref{S1}.
(c) By Gauss' lemma $f(y_1,\dots,Y_i,\dots,y_m)$ is irreducible over $l(y_j:j\neq i)$, hence is the minimal polynomial of $y_i$ up to a unit; it is separable
if and only if its derivative is not zero. Finally $\partial f/\partial Y_i$ has smaller degree in $Y_i$ than $f$, hence is not a multiple of $f$, and so does not
vanish at $(y_1,\dots,y_m)$ unless it is zero.
\end{proof}

\subsection{Levels}
Let $P$ be a point. For $s,t\in P\sm l$ write $s\approx t$ if $s$ is separable algebraic over $l(t)$ and $t$ is separable algebraic over $l(s)$.

\begin{lemma}\label{lem:levels}
The relation $\approx$ is an equivalence relation on $P\sm l$, and for all $s,t\in P\sm l$ there is a unique $j\in\Gamma$ with $s\approx\Frob^j(t)$.
\end{lemma}

\begin{proof}
Reflexivity and symmetry are clear. If $s\approx t$ and $t\approx u$, then $l(s,t,u)$ is separable over $l(s,t)$ and $l(s,t)$ is separable over $l(s)$, so $u$ is separable over $l(s)$;
by symmetry $s\approx u$.
Let $s,t\in P\sm l$ and let $f(S,T)$ be the irreducible relation of Lemma~\ref{lem:sepcrit}. Write $f(S,T)=f_0(S^{p^a},T^{p^b})$ with $a,b\in\Gamma_{\ge0}$ maximal. Then $f_0$ is the irreducible relation
of $(\Frob^as,\Frob^bt)$, and by \ref{S1} and the maximality of $a,b$ both its partial derivatives are non-zero, so $\Frob^as\approx\Frob^bt$ by Lemma~\ref{lem:sepcrit}(c).
Applying the automorphism $\Frob^{-a}$ of $L$, which preserves $l$, gives $s\approx\Frob^{b-a}t$.
If $s\approx\Frob^jt$ and $s\approx\Frob^{j'}t$ with $j<j'$ in $\Gamma$, then $\Frob^jt\approx\Frob^{j'}t$, which is absurd since $\Frob^jt$ is purely inseparable of degree $p^{j'-j}>1$ over $l(\Frob^{j'}t)$.
\end{proof}

\begin{definition}
The set $\Lambda(P)=(P\sm l)/{\approx}$ is the set of \emph{levels} of $P$. The class of $s$ is denoted $\lev(s)$. The group $\Gamma$ acts on $\Lambda(P)$ by
$\lev(s)+j=\lev(\Frob^js)$, and by Lemma~\ref{lem:levels} this makes $\Lambda(P)$ a $\Gamma$-torsor. For $\alpha,\beta\in\Lambda(P)$ we write $\alpha-\beta\in\Gamma$ for the
element with $\alpha=\beta+(\alpha-\beta)$, and $\alpha\ge\beta$ if $\alpha-\beta\geq0$.
\end{definition}

\subsection{Normalized triples}

\begin{definition}
Let $P,Q,R$ be three distinct collinear points. A triple $(x,y,z)$ with $x\in P\sm l$, $y\in Q\sm l$, $z\in R\sm l$ is \emph{normalized}
if each of $x,y,z$ is separable algebraic over the field generated over $l$ by the other two.
\end{definition}

By Lemma~\ref{lem:sepcrit}(c), $(x,y,z)$ is normalized if and only if all three partial derivatives of its irreducible relation are non-zero.

\begin{lemma}[Normalized triples]\label{lem:normalized}
Let $P,Q,R$ be distinct collinear points.
\begin{enumerate}[label=\textup{(\alph*)}]
\item Normalized triples exist.
\item If $(x,y,z)$ is normalized, then $(x',y',z')$ is normalized if and only if there is $j\in\Gamma$ with
$\lev(x')=\lev(x)+j$, $\lev(y')=\lev(y)+j$ and $\lev(z')=\lev(z)+j$.
\end{enumerate}
\end{lemma}

Thus the levels of normalized triples form a single orbit of the diagonal action of $\Gamma$ on $\Lambda(P)\times\Lambda(Q)\times\Lambda(R)$.
We call this orbit, or any of its elements, the \emph{level triple} of $(P,Q,R)$ and denote it by $\lt(P,Q,R)$.

\begin{proof}
(a) Take any $x,y,z$ and their irreducible relation $f$. Write $f=g(X^{p^a},Y^{p^b},Z^{p^c})$ with $a,b,c\in\Gamma_{\ge0}$ maximal. Then $g$ is the irreducible relation of
$(\Frob^ax,\Frob^by,\Frob^cz)$ and by \ref{S1} none of its partial derivatives vanishes.

(b) Being normalized is invariant under $\Frob^j$. It only depends on the levels: if $x'\approx x$, $y'\approx y$, $z'\approx z$ and $(x,y,z)$ is normalized, then
$l(y,z)$ is separable over $l(y',z')$, $x$ is separable over $l(y,z)$ and $x'$ is separable over $l(x)$, so $x'$ is separable over $l(y',z')$; similarly for $y'$ and $z'$.
This proves ``if''.

For ``only if'', let $E\subset L$ be the field of elements separable algebraic over $l(x,y)$. By \ref{S2}, $\Omega_{E|l}$ has basis $dx,dy$.
Since $z\in E$ and the relation $f$ of $(x,y,z)$ has non-vanishing partial derivatives, $f_X\,dx+f_Y\,dy+f_Z\,dz=0$ shows that
\begin{equation}\label{eq:pairwise}
dx,\ dy,\ dz\ \text{ are pairwise linearly independent in }\Omega_{E|l}.
\end{equation}
Let $(x',y',z')$ be normalized. By Lemma~\ref{lem:levels} and the first part of the proof we may assume, after a Frobenius twist, that
$(x',y',z')=(\Frob^ax,\Frob^by,\Frob^cz)$ with $a,b,c\in\Gamma_{\ge0}$ and $\min(a,b,c)=0$; we must show $a=b=c=0$.
All three elements lie in $E$, and each is separable algebraic over the field generated by the other two, so by \ref{S2} its differential lies in the $E$-span of the differentials of the other two.
The differential of $\Frob^ax$ is $dx$ if $a=0$ and $0$ if $a>0$, and similarly for the others.
If exactly one of $a,b,c$ is zero, we get a non-zero differential in the span of two vanishing ones. If exactly two are zero, we get two of the differentials $dx,dy,dz$ each of which is a multiple of the other,
contradicting \eqref{eq:pairwise}. Hence $a=b=c=0$.
\end{proof}

\section{Models and the lax differential map}\label{sec:models}

We keep the assumptions of Section~\ref{sec:levels}.

\begin{definition}
Let $A$ be a transcendence basis of $L|l$. The \emph{model} attached to $A$ is the field
\[ L_A=\{y\in L : y \text{ is separable algebraic over } l(A)\}, \]
and we put $\Omega_A=\Omega_{L_A|l}$.
\end{definition}

\begin{lemma}\label{lem:model}
\begin{enumerate}[label=\textup{(\alph*)}]
\item $\{da:a\in A\}$ is a basis of the $L_A$-vector space $\Omega_A$, and \ref{S3} applies to $E=L_A$.
\item If $j\in\Gamma_{\ge 0}$ and $y\in L$ is separable algebraic over $\Frob^j(L_A)$, then $y\in\Frob^j(L_A)$.
\item $L=\bigcup_{j\in\Gamma_{\ge0}}\Frob^{-j}(L_A)$.
\end{enumerate}
\end{lemma}

\begin{proof}
(a) is \ref{S2}; $l$ is relatively algebraically closed in $L_A$ because it is so in $L$.
(b) Apply $\Frob^{-j}$: $\Frob^{-j}y$ is separable over $L_A$, hence over $l(A)$.
(c) Every $y\in L$ is algebraic over $l(A)$, so $\Frob^jy$ is separable over $l(A)$ for some $j\in\Gamma_{\geq0}$.
\end{proof}

\begin{lemma}[Points in a model]\label{lem:pointsmodel}
Let $P$ be a point. There is a unique level $\lambda_A(P)\in\Lambda(P)$ such that for $y\in P\sm l$:
\begin{enumerate}[label=\textup{(\alph*)}]
\item $y\in\Frob^j(L_A)$ if and only if $\lev(y)\ge\lambda_A(P)+j$ $(j\in\Gamma)$;
\item $y\in L_A$ and $dy\neq0$ in $\Omega_A$ if and only if $\lev(y)=\lambda_A(P)$;
\item the line $L_A\,dy\subset\Omega_A$ is the same for all $y$ as in \textup{(b)}.
\end{enumerate}
\end{lemma}

\begin{proof}
By Lemma~\ref{lem:model}(c), \ref{S3} and since $P$ is stable under $\Frob^{\pm1}$, there is $y_0\in P\cap L_A$, $y_0\notin l$, with $dy_0\neq0$. Put $\lambda_A(P)=\lev(y_0)$.
Let $y\in P\sm l$ and $y\approx\Frob^i(y_0)$ with $i\in\Gamma$. If $i\ge0$, then $y$ is separable over $l(\Frob^iy_0)\subset\Frob^i(L_A)$, so $y\in\Frob^i(L_A)$ by Lemma~\ref{lem:model}(b);
if moreover $i=0$, then $dy\in L_A\,dy_0$, and $dy_0\in L_A\,dy$ because $y_0$ is separable over $l(y)$; if $i>0$ then $dy=0$ by \ref{S2}.
If $i<0$ then $y\notin L_A$: otherwise $y_0$, being separable over $l(\Frob^{-i}y)\subset\Frob^{-i}(L_A)$, would lie in $\Frob^{-i}(L_A)$, and $dy_0=0$.
This proves (b), (c) and (a) for $j=0$; the general case of (a) follows by applying $\Frob^{-j}$.
\end{proof}

\begin{definition}
For a point $P$ we call $\lambda_A(P)$ the \emph{model level} of $P$, an element $y\in P$ with $\lev(y)=\lambda_A(P)$ a \emph{model generator} of $P$, and
\[ \omega_A(P)=L_A\,dy\subset\Omega_A\qquad(y\text{ a model generator of }P) \]
the \emph{line of $P$}. For $\alpha\in\Lambda(P)$ the element $\hgt_A(\alpha)=\alpha-\lambda_A(P)\in\Gamma$ is the \emph{height} of $\alpha$.
\end{definition}

So an element of $P\sm l$ of height $j$ is $\Frob^j$ of a model generator, up to the relation $\approx$.

\begin{lemma}[Lax property]\label{lem:lax}
If $P_1,\dots,P_m$ are distinct points with $\rk\{P_1,\dots,P_m\}<m$, then the lines $\omega_A(P_1),\dots,\omega_A(P_m)$ are linearly dependent.
\end{lemma}

\begin{proof}
Let $y_i$ be a model generator of $P_i$ and let $f\in l[Y_1,\dots,Y_m]$ be a non-zero relation of minimal total degree with $f(y_1,\dots,y_m)=0$. Not all partial derivatives of $f$ are zero,
since otherwise, by \ref{S1}, $f=g^p$ with $p>1$ and $g$ would be a relation of smaller degree. A non-zero partial derivative has smaller degree than $f$, hence does not vanish at $(y_i)$.
Thus $\sum_i f_{Y_i}(y)\,dy_i=0$ is a non-trivial linear relation.
\end{proof}

The lax property is used in the following form, which we call \emph{propagation}: if a point $X$ is collinear with $P,Q$ and also with $R,S$, where
$\omega_A(P)\neq\omega_A(Q)$, $\omega_A(R)\ne\omega_A(S)$ and the planes $\omega_A(P)+\omega_A(Q)$ and $\omega_A(R)+\omega_A(S)$ are distinct, then
$\omega_A(X)$ is the intersection of these two planes.

When $\Gamma\neq0$ the converse of Lemma~\ref{lem:lax} fails: if $a,b\in A$ then $[a]$ and $[a+\Frob(b)]$ are distinct points with the same line $L_A\,da$.
The next lemma describes exactly what happens on a collinear triple.

\begin{lemma}[Trichotomy]\label{lem:trichotomy}
Let $P,Q,R$ be distinct collinear points and let $(i,j,k)\in\Gamma^3$ be the heights of a level triple $(\alpha,\beta,\gamma)$ of $(P,Q,R)$.
\begin{enumerate}[label=\textup{(\alph*)}]
\item If $i<j$, then $k=i$ and $\omega_A(R)=\omega_A(P)$.
\item If $i=j$ and $\omega_A(P)\neq\omega_A(Q)$, then $k=i$, and $\omega_A(R)$ is contained in $\omega_A(P)+\omega_A(Q)$ and different from $\omega_A(P)$ and $\omega_A(Q)$.
\item If $i=j$ and $\omega_A(P)=\omega_A(Q)$, then $k\ge i$.
\item Conversely, if $\omega_A(P),\omega_A(Q),\omega_A(R)$ are pairwise distinct, then $(\lambda_A(P),\lambda_A(Q),\lambda_A(R))$ is a level triple of $(P,Q,R)$, i.e.\ $i=j=k$.
\end{enumerate}
\end{lemma}

\begin{proof}
Shifting $(\alpha,\beta,\gamma)$ diagonally shifts $(i,j,k)$ diagonally, so in (a)--(c) we may assume $i=0\le j$. Let $(x,y,z)$ be a normalized triple with levels $(\alpha,\beta,\gamma)$ and let $f$ be its irreducible relation;
by Lemma~\ref{lem:sepcrit}(c) the elements $f_X,f_Y,f_Z$ of $L$ obtained by evaluating the partial derivatives are non-zero.
Now $x$ is a model generator of $P$ and $y\in L_A$, so $z$, being separable over $l(x,y)$, lies in $L_A$ by Lemma~\ref{lem:model}(b), and by \ref{S2}
\[ f_X\,dx+f_Y\,dy+f_Z\,dz=0\quad\text{in }\Omega_A. \]
(a) Here $dy=0$, so $dz=-(f_X/f_Z)\,dx\ne0$. By Lemma~\ref{lem:pointsmodel}, $z$ is a model generator of $R$ and $\omega_A(R)=\omega_A(P)$.
(b) Here $dx$ and $dy$ are linearly independent and $dz=-(f_X\,dx+f_Y\,dy)/f_Z$ has both coordinates non-zero.
(c) $z\in L_A$ means $k\ge0$.
(d) Let $x,y,z$ be model generators and $f$ their irreducible relation, so $f_X\,dx+f_Y\,dy+f_Z\,dz=0$ with evaluated partial derivatives.
If the polynomial $\partial f/\partial Z$ were zero, then by Lemma~\ref{lem:sepcrit}(b),(c) at least one of $f_X,f_Y$ would be non-zero and $f_X\,dx+f_Y\,dy=0$, contradicting the linear independence of $dx$ and $dy$.
So $z$ is separable over $l(x,y)$; by symmetry $(x,y,z)$ is normalized.
\end{proof}

\section{The quadrangle lemma}\label{sec:quadrangle}

\begin{lemma}[Quadrangle]\label{lem:quadrangle}
Let $A,B,C$ be non-collinear points, let $D$ be a point on $\la A,B\ra$ different from $A,B$, and $E$ a point on $\la B,C\ra$ different from $B,C$.
Let $G$ be a point lying on both $\la D,C\ra$ and $\la A,E\ra$, and let $H$ be a point lying on both $\la A,C\ra$ and $\la B,G\ra$.
Then there are levels $\alpha_X\in\Lambda(X)$ for $X\in\{A,B,C,D,E,G,H\}$ such that for each of the six collinear triples
\[ (A,B,D),\ (B,C,E),\ (D,C,G),\ (A,E,G),\ (A,C,H),\ (B,H,G) \]
the corresponding triple of levels $\alpha_X$ is its level triple.
\end{lemma}

\begin{proof}
First, $G\notin\{A,C,D,E\}$ and $H\notin\{A,B,C,G\}$, by non-collinearity of $A,B,C$; for instance $G=D$ would force $E$ on $\la A,B\ra\cap\la B,C\ra=B$.
Choose a normalized triple $(a,b,d)$ for $(A,B,D)$. By Lemma~\ref{lem:normalized} there are $c\in C$, $e\in E$ such that $(b,c,e)$ is normalized.
Since $a,b,c$ are algebraically independent, they are part of a transcendence basis $\mathcal A$ of $L|l$; we work in the model $L_{\mathcal A}$ and write $\omega=\omega_{\mathcal A}$.
Then $a,b,c$ are model generators and $\omega(A),\omega(B),\omega(C)$ are spanned by $da,db,dc$.
By Lemma~\ref{lem:trichotomy}(b), $d$ and $e$ are model generators and
\[ \omega(D)\ni d_a\,da+d_b\,db,\qquad \omega(E)\ni e_b\,db+e_c\,dc,\qquad d_a,d_b,e_b,e_c\neq0. \]
By propagation, $\omega(G)$ is the intersection of $\omega(D)+\omega(C)$ and $\omega(A)+\omega(E)$, which is spanned by
\[ v=d_a\,da+d_b\,db+\frac{d_b\,e_c}{e_b}\,dc . \]
All three coordinates of $v$ are non-zero, so $\omega(G)$ differs from $\omega(D),\omega(C),\omega(A),\omega(E)$.
By Lemma~\ref{lem:trichotomy}(d) the model levels form level triples of $(D,C,G)$ and of $(A,E,G)$.
Again by propagation, $\omega(H)$ is the intersection of $\omega(A)+\omega(C)$ and $\omega(B)+\omega(G)$, which is spanned by $d_a\,da+\frac{d_be_c}{e_b}dc$;
it differs from $\omega(A),\omega(C),\omega(B),\omega(G)$, and Lemma~\ref{lem:trichotomy}(d) applies to $(A,C,H)$ and $(B,H,G)$.
So $\alpha_X=\lambda_{\mathcal A}(X)$ has the required property.
\end{proof}

\section{Proof of Theorem~\ref{thm:A}}\label{sec:main}

\subsection{Reductions and notation}
The homomorphism $\psi$ extends uniquely to $(K^i)^\times$ by $\psi(x)=p_K^{-j}\psi(x^{p_K^j})$, where $p_K$ is the characteristic exponent of $K$; the extension satisfies \ref{A1}--\ref{A4} with $k$ replaced by its perfect closure $k^i\subset K^i$,
and $k^i\cap K=k$ by \ref{A1}. Also $V_{L|l}=V_{L^i|l^i}$. So we may and do assume that $K$ and $L$ are perfect; then $k$ is perfect.
Note that at this stage $K$ and $L$ may have different characteristics; until this is excluded (Proposition~\ref{prop:tau}), $\Frob$ always denotes the Frobenius of $L$ and we make no use of the Frobenius of $K$.
We write $p$ for the characteristic exponent of $L$, $\Gamma=\Gamma(L)$, and $V=V_{L|l}$. By \ref{A1}, $\psi(-1)=0$ and $\psi(cx)=\psi(x)$ for $c\in k^\times$.

For $v\in V\sm0$ with representative $(N,a)$ we put $\lev(v)=\lev(a)\in\Lambda(Q_v)$. This is well defined: if $(N',a')$ is another representative then $a^{N'}=ca'^N$, and $a\approx a^{N'}$, $a'\approx a'^N$ because $N,N'$ are prime to $p$.
Clearly $\lev(pv)=\lev(v)+1$ if $p>1$, and $\lev(-v)=\lev(v)$.

We shall use the following consequence of \ref{A1} and \ref{A2} several times: for $x,y\in K^\times$ with $x+y\neq0$ and $x/y\notin k$,
\begin{equation}\label{eq:sum}
\psi(x+y)=\psi(y)+\psi\bigl(1-(-x/y)\bigr),\qquad\text{and}\quad \psi\bigl(1-(-x/y)\bigr)\in G_{Q_{x/y}}\sm0 .
\end{equation}
In particular, if $x,y$ are $\psi$-independent then $\psi(x+y)\ne0$ and $Q_x,Q_y,Q_{x+y}$ are three distinct collinear points. Indeed $Q_{x/y}$ lies on $\la Q_x,Q_y\ra$ and is different from $Q_y$ (otherwise $\psi(x)=\psi(x/y)+\psi(y)\in G_{Q_y}$), so by \eqref{eq:sum} the element $\psi(x+y)$ is the sum of a non-zero element of $G_{Q_y}$ and a non-zero element of $G_{Q_{x/y}}$: it is non-zero, its point lies on $\la Q_x,Q_y\ra$ and is not $Q_y$; by symmetry it is not $Q_x$.
For every $x\in K\sm k$ we have $1-x\notin k$, and $Q_{1-x}=Q_x$ by \ref{A2}.

\subsection{Step A: the levels of \texorpdfstring{$\psi(x)$}{psi(x)} and \texorpdfstring{$\psi(1-x)$}{psi(1-x)}}\label{sec:stepA}

\begin{proposition}\label{prop:levels}
For every $x\in K\sm k$ we have $\lev\psi(1-x)=\lev\psi(x)$.
\end{proposition}

If $\Gamma=0$ there is nothing to prove. In this subsection we assume $\Gamma=\ZZ$, so $d=3$, and by \ref{A4} the field $K$ has positive characteristic; as $K$ is perfect, $K^\times$ is divisible by its characteristic, and therefore
\begin{equation}\label{eq:nohom}
\text{every homomorphism }K^\times\to\ZZ\text{ is trivial.}
\end{equation}
For $s\in K\sm k$ put
\[ n(s)=\lev\psi(1-s)-\lev\psi(s)\in\ZZ . \]
Since $\psi(1-1/s)=\psi(1-s)-\psi(s)$ and $\lev\psi(1/s)=\lev\psi(s)$, we have $n(1/s)=\lev\bigl(\psi(1-s)-\psi(s)\bigr)-\lev\psi(s)$. Obviously $n(1-s)=-n(s)$.

For $\psi$-independent $u,w$ the points $Q_u,Q_w,Q_{u+w}$ are distinct and collinear, and we define $c(u,w)\in\ZZ$ by the condition that
\begin{equation}\label{eq:defc}
\bigl(\lev\psi(u),\ \lev\psi(w)+c(u,w),\ \ast\bigr)\ \text{ is the level triple of }(Q_u,Q_w,Q_{u+w}).
\end{equation}
Clearly $c(w,u)=-c(u,w)$.

\begin{lemma}[The gauge of a sum]\label{lem:gauge}
Let $u,w$ be $\psi$-independent and $s=-u/w$. Then
\[ c(u,w)=\begin{cases}-n(s)&\text{if }n(s)>0,\\ n(1/s)&\text{if }n(s)=0\text{ and }n(1/s)>0,\\ 0&\text{if }n(s)<0,\text{ or }n(s)=n(1/s)=0,\end{cases} \]
and these cases are exhaustive. In particular $c(u,w)$ depends only on $s$; we denote it by $\nu(s)$.
\end{lemma}

\begin{proof}
Choose representatives with a common $N$: $N\psi(s)=[a]$, $N\psi(w)=[y]$, $N\psi(1-s)=[e]$. Then $a,e\in Q_s$, the elements $a,y$ are algebraically independent, and by \eqref{eq:sum}
\[ N\psi(u)=[ay],\qquad N\psi(u+w)=[ye],\qquad \lev(e)=\lev(a)+n(s),\qquad \lev(e/a)=\lev(a)+n(1/s). \]
In each case we exhibit a normalized triple for $(Q_u,Q_w,Q_{u+w})=([ay],[y],[ye])$. To check that $Y$ is separable over $l(X,Z)$ we use the following remark: if $X,Y,Z$ lie in a field $E'$ as in \ref{S2} and
$f$ is their irreducible relation, then $f_X\,dX+f_Y\,dY+f_Z\,dZ=0$, and by Lemma~\ref{lem:sepcrit} not all of $f_X,f_Y,f_Z$ vanish.

\emph{Case $i=n(s)>0$.} Write $e=\Frob^i(e_0)$ with $e_0\approx a$, and put $A=\Frob^i(a)$. We claim that $(X,Y,Z)=(Ay^{p^i},\,y,\,ye)$ is normalized.
Indeed $e$ is separable over $l(A)$, so $Z$ is separable over $l(X,Y)=l(A,y)$; and $A$ is separable over $l(e)$, so $X$ is separable over $l(Y,Z)=l(y,e)$.
Let $E'$ be the field of elements separable over $l(A,y)$, with basis $dA,dy$ of $\Omega_{E'|l}$. Then $dX=y^{p^i}dA$ and $dZ=e\,dy+y\,e_A\,dA$, where $de=e_A\,dA$, so the coefficient of $dy$ in the relation above is $f_Y+f_Ze$ and that of $dA$ is $f_Xy^{p^i}+f_Zye_A$.
If $f_Y=0$ we get $f_Z=0$ and then $f_X=0$, a contradiction. So $Y$ is separable over $l(X,Z)$.
The levels of this triple are $(\lev(ay)+i,\lev(y),\lev(ye))$, hence $c(u,w)=-i$.

\emph{Case $i=n(s)<0$.} Put $j=-i$, $\alpha=\Frob^{-j}(a)$, $\eta=\Frob^{-j}(y)$, so $e\approx\alpha$. We claim that $(X,Y,Z)=(\alpha\eta,\,\eta,\,\eta^{p^j}e)$ is normalized; note $Z=ye$.
As before $Z$ is separable over $l(X,Y)=l(\alpha,\eta)$ and $X$ is separable over $l(Y,Z)=l(\eta,e)$. In the field $E'$ of elements separable over $l(\alpha,\eta)$ we have $dX=\eta\,d\alpha+\alpha\,d\eta$ and $dZ=\eta^{p^j}e_\alpha\,d\alpha$, where $de=e_\alpha\,d\alpha\neq0$ because $e\approx\alpha$ (Lemma~\ref{lem:pointsmodel}(b)).
The coefficient of $d\eta$ is $f_X\alpha+f_Y$ and that of $d\alpha$ is $f_X\eta+f_Z\eta^{p^j}e_\alpha$; if $f_Y=0$ then $f_X=0$ and $f_Z=0$. So the triple is normalized.
Its levels are $(\lev(ay)-j,\lev(y)-j,\lev(ye))$, hence $c(u,w)=0$.

\emph{Case $n(s)=0$.} Then $e\approx a$, so $e/a$ is separable over $l(a)$ and is not constant; hence $\lev(e/a)\ge\lev(a)$, i.e.\ $n(1/s)\geq0$.
(In general, if $g\in Q_s\sm l$ is separable over $l(a)$ then $\lev(g)\ge\lev(a)$: if $a\approx\Frob^j(g)$ with $j>0$, then $g$ would be separable over $l(\Frob^jg)$, while it is purely inseparable of degree $p^j$ over this field.)
Suppose first $n(1/s)=0$. We claim that $(X,Y,Z)=(ay,\,y,\,ye)$ is normalized. The first two separability conditions are clear. In the field $E'$ of elements separable over $l(a,y)$,
$dX=y\,da+a\,dy$ and $dZ=e\,dy+y\,e_a\,da$; the coefficient of $dy$ is $f_Xa+f_Y+f_Ze$ and that of $da$ is $f_Xy+f_Zye_a$. If $f_Y=0$ this gives $f_X=-f_Ze_a$ and $f_Z(e-ae_a)=0$.
Now $e-ae_a=-a^2(e/a)_a$, where $d(e/a)=(e/a)_a\,da$, and $d(e/a)\ne0$ because $\lev(e/a)=\lev(a)$ (Lemma~\ref{lem:pointsmodel}(b), for a transcendence basis containing $a$ and $y$). So $f_Z=f_X=0$, a contradiction. Hence the triple is normalized and $c(u,w)=0$.

Suppose finally $n(s)=0$ and $n(1/s)>0$. Apply the first case to the pair $(w,u)$, for which $-w/u=1/s$: it gives $c(w,u)=-n(1/s)$, so $c(u,w)=n(1/s)$.
\end{proof}

\begin{lemma}[Cocycle]\label{lem:cocycle}
If $u,w,z\in K$ are $\psi$-independent, then $c(u,z)=c(u,w)+c(w,z)$.
\end{lemma}

\begin{proof}
Apply Lemma~\ref{lem:quadrangle} to the points
\[ A=Q_u,\quad B=Q_w,\quad C=Q_z,\quad D=Q_{u+w},\quad E=Q_{w+z},\quad G=Q_{u+w+z},\quad H=Q_{u+z}. \]
The points $A,B,C$ are not collinear; $D$ lies on $\la A,B\ra$ and $E$ on $\la B,C\ra$, and they differ from $A,B$ resp.\ $B,C$ by the remark after \eqref{eq:sum}.
By \eqref{eq:sum}, applied to $(u+w)+z$ and to $u+(w+z)$, the point $G$ lies on $\la D,C\ra$ and on $\la A,E\ra$; in particular $G\neq B$, since $\la D,C\ra\cap\la A,B\ra=D$. Applying \eqref{eq:sum} to $u+z$ and to $(u+w+z)+(-w)$ shows that $H$ lies on $\la A,C\ra$ and on $\la B,G\ra$.
So there are levels $\alpha_A,\alpha_B,\alpha_C$ which are the first two entries of level triples of $(A,B,D)$, of $(B,C,E)$ and of $(A,C,H)$. By \eqref{eq:defc} this means
\[ \alpha_B-\lev\psi(w)=\alpha_A-\lev\psi(u)+c(u,w),\qquad \alpha_C-\lev\psi(z)=\alpha_B-\lev\psi(w)+c(w,z), \]
and $\alpha_C-\lev\psi(z)=\alpha_A-\lev\psi(u)+c(u,z)$. The claim follows.
\end{proof}

\begin{proof}[Proof of Proposition~\ref{prop:levels}]
Let $\nu$ be the function of Lemma~\ref{lem:gauge}. Let $s,s'\in K$ be $\psi$-independent. By \ref{A3} there is $z$ such that $s,s',z$ are $\psi$-independent (one of $x_1,x_2,x_3$ will do). Put $w=-s'z$ and $u=ss'z$.
Then $u,w,z$ are $\psi$-independent, $-u/w=s$, $-w/z=s'$ and $-u/z=-ss'$, so Lemma~\ref{lem:cocycle} gives
\begin{equation}\label{eq:nu}
\nu(-ss')=\nu(s)+\nu(s')\qquad\text{for all $\psi$-independent }s,s'.
\end{equation}
Applying this to $(-s,s')$ and to $(s,-s')$ we get $\nu(-s)-\nu(s)=\nu(-s')-\nu(s')$. So $\nu(-s)-\nu(s)$ takes the same value at any two $\psi$-independent elements, hence is a constant $\delta_0$.
Then $\nu(ss')=\nu(s)+\nu(s')+\delta_0$ by \eqref{eq:nu} applied to $(-s,s')$, and so $\nu(-ss')=\nu(s)+\nu(s')+2\delta_0$; comparing with \eqref{eq:nu} gives $\delta_0=0$. Thus
\[ \nu(ss')=\nu(s)+\nu(s')\qquad\text{for all $\psi$-independent }s,s'. \]
For $c\in k^\times$ and $w\notin k$ the integer $\varrho(c)=\nu(cw)-\nu(w)$ does not depend on $w$: if $w,w'$ are $\psi$-independent, then $\nu(cw)+\nu(w')=\nu(cww')=\nu(w)+\nu(cw')$, and two elements with the same point are compared with a third one.
Define $\chi\colon K^\times\to\ZZ$ by $\chi=\nu$ on $K\sm k$ and $\chi=\varrho$ on $k^\times$. We claim that $\chi$ is a homomorphism. This is clear for pairs of constants, for a constant and a non-constant, and for $\psi$-independent pairs.
Let $x,y\notin k$ with $Q_x=Q_y$ and choose $w$ with $Q_w\neq Q_x$. If $xy\notin k$, then $\nu(xy)+\nu(w)=\nu(xyw)=\nu(x)+\nu(yw)=\nu(x)+\nu(y)+\nu(w)$. If $xy=c\in k$, then $\varrho(c)+\nu(w)=\nu(cw)=\nu(x)+\nu(yw)=\nu(x)+\nu(y)+\nu(w)$.
By \eqref{eq:nohom}, $\chi=0$, so $\nu=0$.

By Lemma~\ref{lem:gauge}, $\nu(s)=0$ excludes $n(s)>0$. Hence $n(s)\le0$ for all $s\in K\sm k$, and since $n(1-s)=-n(s)$ we conclude $n=0$.
\end{proof}

\subsection{The adapted model}\label{ss:model}
Let $x_1,\dots,x_d$ be as in \ref{A3}, choose representatives $(N_i,a_i)$ of $\psi(x_i)$, and extend $a_1,\dots,a_d$ to a transcendence basis $A$ of $L|l$. Put
\[ F=L_A,\qquad \Omega=\Omega_A=\Omega_{L_A|l}, \]
and for $v\in V\sm0$ let $\hgt(v)=\hgt_A(\lev v)\in\Gamma$ be its height. (If $\Gamma=0$ then $F=L$, $\Omega=\Omega_{L|l}$ and all heights are $0$.) Let
\[ W=\{0\}\cup\{v\in V\sm0:\hgt(v)\geq0\}. \]

\begin{lemma}\label{lem:W}
\begin{enumerate}[label=\textup{(\alph*)}]
\item $W$ is a subgroup of $V$, and there is a homomorphism $\dlog_A\colon W\to\Omega$ such that $\dlog_A(v)=N^{-1}\,da/a$ for every representative $(N,a)$ of $v\in W\sm0$.
\item For $v\in W\sm0$: $\dlog_A(v)\neq0$ if and only if $\hgt(v)=0$, and then $F\cdot\dlog_A(v)=\omega_A(Q_v)$.
\item If $\Gamma=0$, then $W=V$ and $\dlog_A$ is injective. If $\Gamma=\ZZ$, then $\ker(\dlog_A)=pW$, $V=\bigcup_jp^{-j}W$ and $\bigcap_jp^jW=0$.
\end{enumerate}
\end{lemma}

\begin{proof}
Let $v\in V\sm0$ have representative $(N,a)$. By Lemma~\ref{lem:pointsmodel}, $a\in L_A$ if and only if $\hgt(v)\ge0$, and then $da\ne0$ if and only if $\hgt(v)=0$, in which case $F\,da=\omega_A(Q_v)$.
If $(N',a')$ is another representative, $a^{N'}=ca'^N$ gives $N'\,da/a=N\,da'/a'$; so $\dlog_A$ is well defined, and (b) holds.
Given $v,v'\in W$ choose representatives with a common $N$; then $N(v+v')=[aa']$ with $aa'\in L_A$, so $v+v'\in W$ and $\dlog_A$ is additive. This proves (a).
(c) follows from $\hgt(pv)=\hgt(v)+1$ and \ref{V1}: if $\Gamma=0$ and $\dlog_A(v)=0$ then $da=0$, so $a$ is algebraic over $l$ by \ref{S3}, i.e.\ $v=0$.
\end{proof}

Put
\[ K_0=\{0\}\cup\psi^{-1}(W),\qquad \Delta=\dlog_A\circ\psi\colon K_0^\times\to\Omega . \]
An element $x\in K_0^\times$ is \emph{primitive} if $\Delta(x)\ne0$, i.e.\ $x\notin k$ and $\hgt\psi(x)=0$. Two primitive elements $x,y$ are \emph{transverse} if $\Delta(x),\Delta(y)$ are linearly independent; by Lemma~\ref{lem:W}(b) this implies $Q_x\neq Q_y$, so transverse elements are $\psi$-independent.

\begin{lemma}\label{lem:K0}
\begin{enumerate}[label=\textup{(\alph*)}]
\item $K_0$ is a subfield of $K$ containing $k$. If $\Gamma=0$, then $K_0=K$ and every $x\in K\sm k$ is primitive.
\item If $x\in K_0\sm k$ then $1-x\in K_0\sm k$, and $1-x$ is primitive if and only if $x$ is. If $x$ is primitive, then
\[ \Delta(1-x)=\mu(x)\,\Delta(x)\qquad\text{for a unique }\mu(x)\in F^\times,\text{ and }\mu(x)\in Q_x . \]
\item $x_1,\dots,x_d$ are primitive and $\Delta(x_1),\dots,\Delta(x_d)$ are linearly independent. For any three lines in $\Omega$ there is a primitive $z$ such that $F\Delta(z)$ is none of them.
\end{enumerate}
\end{lemma}

\begin{proof}
(b) By Proposition~\ref{prop:levels}, $\psi(1-x)$ and $\psi(x)$ have the same height, and they lie in the same point $Q_x$. So $1-x\in K_0$, it is primitive if and only if $x$ is, and for primitive $x$ both $\Delta(x)$ and $\Delta(1-x)$ span $\omega_A(Q_x)$.
With representatives $(N,a)$, $(N,b)$ of $\psi(x),\psi(1-x)$ we have $b\approx a$, so $db=b_a\,da$ with $b_a$ in the field of elements separable over $l(a)$, which is contained in $Q_x$; thus $\mu(x)=ab_a/b\in Q_x$.

(a) $K_0^\times$ is a subgroup of $K^\times$ containing $k^\times$. Let $x,y\in K_0^\times$ with $x+y\ne0$. If $x/y\in k$ then $x+y\in k^\times y$. Otherwise $-x/y\in K_0\sm k$, so $1+x/y\in K_0$ by (b) and $x+y=y(1+x/y)\in K_0$.
The statement for $\Gamma=0$ is Lemma~\ref{lem:W}(c).

(c) The $a_i$ belong to $A$, so $\hgt\psi(x_i)=0$ and $\Delta(x_i)=N_i^{-1}da_i/a_i$. A monomial $z=x_1^{i_1}\cdots x_d^{i_d}$ has $\Delta(z)=\sum i_j\Delta(x_j)$.
If $\Gamma=0$, the elements $x_1x_2^n$, $n\ge0$, give infinitely many lines. If $\Gamma=\ZZ$, then $d=3$ and the monomials give the $p^2+p+1\ge7$ lines $F(\sum i_j\Delta(x_j))$, $(i_1:i_2:i_3)\in\mathbf P^2(\FF_p)$.
\end{proof}

By Lemma~\ref{lem:K0}(c), any two primitive elements are transverse to a common third one.

\subsection{Step B: the functions \texorpdfstring{$\mu$}{mu} and \texorpdfstring{$t$}{t}}

\begin{lemma}[One variable]\label{lem:onevar}
Let $x$ be primitive. Then $1-x$ and $1/x$ are primitive,
\[ \mu(1-x)=\mu(x)^{-1},\qquad \mu(1/x)=1-\mu(x),\qquad \mu(x)\neq1, \]
and the element
\[ t(x)=\frac{\mu(x)}{\mu(x)-1}\in F\sm\{0,1\} \]
satisfies $t(1-x)=1-t(x)$, $t(1/x)=1/t(x)$ and $t(x)\in Q_x$.
\end{lemma}

\begin{proof}
$1-x$ is primitive by Lemma~\ref{lem:K0}(b), and $1/x$ is primitive because $\Delta(1/x)=-\Delta(x)$. From $\Delta(x)=\Delta(1-(1-x))=\mu(1-x)\mu(x)\Delta(x)$ we get the first formula.
Since $1-1/x=-(1-x)/x$ and $\Delta(-1)=0$,
\[ \mu(1/x)\Delta(1/x)=\Delta(1-1/x)=\Delta(1-x)-\Delta(x)=(\mu(x)-1)\Delta(x), \]
which is the second formula; and $\mu(1/x)\ne0$ gives $\mu(x)\ne1$. The identities for $t$ are a direct computation: $t(1-x)=\mu^{-1}/(\mu^{-1}-1)=1/(1-\mu)=1-t(x)$ and $t(1/x)=(1-\mu)/(-\mu)=1/t(x)$, with $\mu=\mu(x)$.
\end{proof}

\begin{lemma}[Two variables]\label{lem:twovar}
If $x,y$ are transverse, then $xy$ is primitive and $t(xy)=t(x)\,t(y)$.
\end{lemma}

\begin{proof}
$\Delta(xy)=\Delta(x)+\Delta(y)\ne0$. Let $u=x(1-y)/(x-1)\in K_0$. Then
\[ \Delta(u)=\Delta(x)+\Delta(1-y)-\Delta(1-x)=(1-\mu(x))\Delta(x)+\mu(y)\Delta(y)\neq0, \]
so $u$ is primitive. Since $1-xy=(1-x)(1-u)$,
\[ \mu(xy)\bigl(\Delta(x)+\Delta(y)\bigr)=\Delta(1-xy)=\mu(x)\Delta(x)+\mu(u)\bigl((1-\mu(x))\Delta(x)+\mu(y)\Delta(y)\bigr). \]
Comparing coefficients, $\mu(xy)=\mu(x)+\mu(u)(1-\mu(x))$ and $\mu(xy)=\mu(u)\mu(y)$. Eliminating $\mu(u)$ gives $\mu(xy)\bigl(\mu(x)+\mu(y)-1\bigr)=\mu(x)\mu(y)$, that is
\[ 1-\frac1{\mu(xy)}=\Bigl(1-\frac1{\mu(x)}\Bigr)\Bigl(1-\frac1{\mu(y)}\Bigr). \]
As $1-1/\mu=1/t$, this is $t(xy)=t(x)t(y)$.
\end{proof}

\subsection{The field embedding}

Let
\[ \mathcal N=\{x\in K_0^\times:\Delta(x)=0\}, \]
a subgroup of $K_0^\times$ containing $k^\times$; so $K_0^\times$ is the disjoint union of $\mathcal N$ and the set of primitive elements. If $x\in\mathcal N$ and $z$ is primitive, then $xz$ is primitive and $\Delta(xz)=\Delta(z)$.

\begin{lemma}\label{lem:tconst}
\begin{enumerate}[label=\textup{(\alph*)}]
\item There is a homomorphism $\tau\colon\mathcal N\to F^\times$ such that $t(xz)=\tau(x)\,t(z)$ for all $x\in\mathcal N$ and all primitive $z$.
\item $\tau(-1)=-1$, and for transverse $x,y$ the element $x+y$ is primitive, transverse to $x$ and to $y$, and
\[ t(x+y)=t(x)+t(y). \]
\end{enumerate}
\end{lemma}

\begin{proof}
(a) Let $x\in\mathcal N$. For transverse $z,z'$ the pairs $(xz,z')$ and $(xz',z)$ are transverse, so Lemma~\ref{lem:twovar} gives $t(xz)t(z')=t(xzz')=t(xz')t(z)$, i.e.\ $t(xz)/t(z)=t(xz')/t(z')$.
Since any two primitive elements are transverse to a common third one, the quotient $t(xz)/t(z)$ is the same for all primitive $z$, and we call it $\tau(x)$.
For $x,x'\in\mathcal N$ we have
\[ \tau(xx')=\frac{t(x\cdot x'z)}{t(x'z)}\cdot\frac{t(x'z)}{t(z)}=\tau(x)\tau(x'), \]
because $x'z$ is primitive.

(b) Let $x,y$ be transverse. Then $-x/y$ is primitive and so is $1+x/y$, with $\Delta(1+x/y)=\mu(-x/y)(\Delta(x)-\Delta(y))$. Hence $x+y=y(1+x/y)$ lies in $K_0$ and
$\Delta(x+y)=\mu(-x/y)\Delta(x)+(1-\mu(-x/y))\Delta(y)$; both coefficients are non-zero by Lemma~\ref{lem:onevar}, so $x+y$ is primitive and transverse to $x$ and to $y$.
The pair $(y,1+x/y)$ is transverse, and so is $(-x,1/y)$. By Lemmas~\ref{lem:onevar} and \ref{lem:twovar} and (a),
\[ t(x+y)=t(y)\bigl(1-t(-x/y)\bigr)=t(y)\Bigl(1-\tau(-1)\frac{t(x)}{t(y)}\Bigr)=t(y)-\tau(-1)\,t(x). \]
By symmetry this equals $t(x)-\tau(-1)t(y)$, so $(1+\tau(-1))(t(x)-t(y))=0$. If $\tau(-1)\ne-1$, then $t(x)=t(y)$ for all transverse pairs, so $t$ is constant on primitive elements; by Lemma~\ref{lem:twovar} the constant is $1$, which is not a value of $t$.
Hence $\tau(-1)=-1$ and $t(x+y)=t(x)+t(y)$.
\end{proof}

We extend $\tau$ to a map $\tau\colon K_0\to F$ by $\tau(0)=0$ and $\tau(x)=t(x)$ for primitive $x$. For $x\in K_0^\times$ let $\ell(x)=F\Delta(x)$; this is a line in $\Omega$ if $x$ is primitive, and $0$ if $x\in\mathcal N$.
By Lemma~\ref{lem:twovar} and Lemma~\ref{lem:tconst}(a),
\begin{equation}\label{eq:P3}
t(xz)=\tau(x)\,t(z)\qquad\text{for all $x\in K_0^\times$ and all primitive $z$ with }F\Delta(z)\neq\ell(x),
\end{equation}
and in this situation $xz$ is primitive.

\begin{proposition}\label{prop:tau}
$\tau\colon K_0\to F$ is a field embedding. In particular $K$ and $L$ have the same characteristic.
\end{proposition}

\begin{proof}
\emph{(1) Multiplicativity.} Let $x,y\in K_0^\times$. If $x,y$ are primitive and transverse, $\tau(xy)=\tau(x)\tau(y)$ is Lemma~\ref{lem:twovar}. Assume this is not the case; then $\ell(y)\subset\ell(x)$ or $\ell(x)\subset\ell(y)$, and by symmetry we may assume $\ell(y)\subset\ell(x)$.
By Lemma~\ref{lem:K0}(c) there is a primitive $z$ such that $F\Delta(z)$ is different from $\ell(x)$, $\ell(y)$ and $\ell(xy)$.
By \eqref{eq:P3}, $t(xyz)=\tau(xy)t(z)$, and $yz$ is primitive with $t(yz)=\tau(y)t(z)$. Moreover $\Delta(yz)=\Delta(y)+\Delta(z)\notin\ell(x)$, because $\Delta(y)\in\ell(x)$ and $\Delta(z)\notin\ell(x)$.
So \eqref{eq:P3} applies to $x$ and $yz$ and gives $t(xyz)=\tau(x)t(yz)=\tau(x)\tau(y)t(z)$. Hence $\tau(xy)=\tau(x)\tau(y)$.

\emph{(2)} $\tau(1-s)=1-\tau(s)$ for all $s\in K_0\sm\{0,1\}$. For primitive $s$ this is Lemma~\ref{lem:onevar}. Let $s\in\mathcal N$; then $1-s\in\mathcal N$ (by Lemma~\ref{lem:K0}(b) if $s\notin k$).
Choose transverse $x,y$. Then $(1-s)x$ is primitive with $\Delta((1-s)x)=\Delta(x)$, so it is transverse to $y$, and by Lemma~\ref{lem:tconst}
\[ t\bigl((1-s)x+y\bigr)=\tau(1-s)\,t(x)+t(y). \]
On the other hand $x+y$ is primitive and transverse to $x$, hence to $-sx$, and $t(-sx)=\tau(-s)t(x)=-\tau(s)t(x)$. So
\[ t\bigl((x+y)+(-sx)\bigr)=t(x)+t(y)-\tau(s)\,t(x). \]
Comparing, $\tau(1-s)=1-\tau(s)$.

\emph{(3) Additivity.} Let $x,y\in K_0^\times$ with $x+y\neq0$. By (1), (2) and $\tau(-1)=-1$,
\[ \tau(x+y)=\tau(y)\,\tau\bigl(1-(-x/y)\bigr)=\tau(y)\bigl(1+\tau(x)/\tau(y)\bigr)=\tau(x)+\tau(y). \]
Also $\tau(-x)=-\tau(x)$ and $\tau(1)=1$. So $\tau$ is a ring homomorphism from the field $K_0$ to $F$, hence injective, and $\chr K=\chr K_0=\chr F=\chr L$.
\end{proof}

From now on $K$ and $L$ have the same characteristic exponent $p$, and $\Frob$ denotes the Frobenius of either field; $\tau$ commutes with $\Frob$.
For $x\in K\sm k$ we have $\psi(\Frob x)=p\,\psi(x)$, which has height $\hgt\psi(x)+1$ if $p>1$. Since $K$ is perfect, this shows:
\begin{equation}\label{eq:roots}
\begin{minipage}{0.86\textwidth}
$\Frob(K_0)\subset K_0$ and $K=\bigcup_{j\in\Gamma_{\ge0}}\Frob^{-j}(K_0)$; every $x\in K_0\sm k$ can be written uniquely as $x=\Frob^j(x_0)$ with $x_0$ primitive and $j\in\Gamma_{\ge0}$, namely $j=\hgt\psi(x)$.
\end{minipage}
\end{equation}
We call $x_0$ the \emph{primitive root} of $x$; then $\tau(x)=\Frob^j(t(x_0))$.

\subsection{Step C: identification}

\begin{lemma}\label{lem:tauinQ}
$\tau(x)\in Q_x\sm l$ for every $x\in K_0\sm k$, and $\tau(k)\subset l$.
\end{lemma}

\begin{proof}
For primitive $x$ we have $\tau(x)=t(x)\in Q_x$ by Lemma~\ref{lem:onevar}; by \eqref{eq:roots}, $\tau(x)=\Frob^j(t(x_0))\in Q_{x_0}=Q_x$ for all $x\in K_0\sm k$.
For $c\in k^\times$ and primitive $x$, $\tau(c)=t(cx)/t(x)\in Q_x$; taking two transverse elements $x$ gives $\tau(c)\in l$.

Suppose $\tau(x)\in l$ for some $x\in K_0\sm k$; replacing $x$ by its primitive root we may assume that $x$ is primitive, and we put $c=\tau(x)=t(x)\in l\sm\{0,1\}$.
By the definition of $t$, $\Delta(1-x)=\kappa\,\Delta(x)$ with $\kappa=\mu(x)=c/(c-1)\in l$. Let $(N,a)$, $(N,b)$ be representatives of $\psi(x)$, $\psi(1-x)$. They are model generators of $Q_x$, so $b\approx a$, and $db/b=\kappa\,da/a$ in $\Omega$.
Let $E$ be the field of elements of $L$ separable algebraic over $l(a)$. Then $E\subset F$ and $\Omega_{E|l}=E\,da$ by \ref{S2}, and $\Omega_{E|l}\to\Omega$ is injective because $da\neq0$ in $\Omega$.
Let $\delta$ be the derivation of $E$ with $df=\delta(f)\,da/a$. Then $\delta(l)=0$, $\delta(a)=a$ and $\delta(b)=\kappa b$, so $\kappa\in\FF$ by Lemma~\ref{lem:R}.
Hence $c=\kappa/(\kappa-1)\in\FF$. But $\tau$ is a field embedding, so it maps the prime field of $K$ onto $\FF$, and $\tau(x)=c$ forces $x$ to lie in the prime field of $K$, contradicting $x\notin k$.
\end{proof}

Recall that $K$ and $L$ have the same characteristic by Proposition~\ref{prop:tau}.

\begin{lemma}\label{lem:theta}
There is $\theta\in\FF^\times$ such that for every primitive $x$
\[ d\tau(x)\neq0\qquad\text{and}\qquad \Delta(x)=\theta\,\dlog\tau(x). \]
\end{lemma}

\begin{proof}
Let $x$ be primitive. By Lemma~\ref{lem:tauinQ}, $\tau(x)\in(Q_x\cap F)\sm l$, so its height $e(x)$ is $\ge0$.
Let $x,y$ be transverse and $(N,a)$, $(N,b)$ representatives of $\psi(x),\psi(y)$. Then $(N,ab)$ represents $\psi(xy)$, all three have height $0$, and $(a,b,ab)$ is a normalized triple of $(Q_x,Q_y,Q_{xy})$.
So is $(\tau x,\tau y,\tau(xy))$, because $\tau(xy)=\tau(x)\tau(y)$ and $\tau x,\tau y$ are algebraically independent. By Lemma~\ref{lem:normalized}, $e(x)=e(y)$. Hence $e$ is a constant $e_0$ on primitive elements, and
$\sigma_x=\Frob^{-e_0}\tau(x)$ is a model generator of $Q_x$.

So $\Delta(x)=\theta(x)\,\dlog\sigma_x$ with $\theta(x)\in F^\times$. For transverse $x,y$ we have $\sigma_{xy}=\sigma_x\sigma_y$, so
$\theta(xy)(\dlog\sigma_x+\dlog\sigma_y)=\theta(x)\dlog\sigma_x+\theta(y)\dlog\sigma_y$ and $\theta(x)=\theta(y)$. Thus $\theta$ is constant.
With a representative $(N,a)$ of $\psi(x)$ we have $a\approx\sigma_x$ and $\theta=N^{-1}(da/a)/(d\sigma_x/\sigma_x)\in Q_x$; taking two transverse $x$ shows $\theta\in l$.
As in the proof of Lemma~\ref{lem:tauinQ}, let $\delta$ be the derivation of the field of elements separable algebraic over $l(\sigma_x)$ with $df=\delta(f)\,d\sigma_x/\sigma_x$. Then $\delta(\sigma_x)=\sigma_x$ and $\delta(a)=N\theta\,a$, so $N\theta\in\FF$ by Lemma~\ref{lem:R}, and $\theta\in\FF^\times$.

Finally $\sigma_{1-x}=1-\sigma_x$, so
\[ \mu(x)=\frac{\Delta(1-x)}{\Delta(x)}=\frac{\dlog(1-\sigma_x)}{\dlog\sigma_x}=\frac{\sigma_x}{\sigma_x-1},\qquad\text{hence}\quad \tau(x)=t(x)=\sigma_x=\Frob^{-e_0}\tau(x). \]
As $\tau(x)\notin l$ this forces $e_0=0$.
\end{proof}

\begin{proposition}\label{prop:rational}
There is $r\in\QQ^\times$ with $\psi(x)=r\,[\tau(x)]$ for all $x\in K_0^\times$.
\end{proposition}

\begin{proof}
Note that $x\mapsto[\tau(x)]$ is a homomorphism $K_0^\times\to W$ which vanishes on $k^\times$ and satisfies $\dlog_A[\tau(x)]=\dlog\tau(x)$.

\emph{The case $\Gamma=0$.} Then $\theta=m/n$ with integers $m,n\ne0$, and $\dlog_A\bigl(n\psi(x)-m[\tau(x)]\bigr)=0$ for all $x\in K^\times$ by Lemma~\ref{lem:theta}. Since $\dlog_A$ is injective, $\psi=\frac mn[\tau]$.

\emph{The case $\Gamma=\ZZ$.} We first show:
\begin{equation}\label{eq:lift}
\begin{minipage}{0.86\textwidth}
if $\psi'\colon K_0^\times\to W$ is a homomorphism with $\psi'(k^\times)=0$ and $\psi'(x)\in G_{Q_x}$ for all $x\in K_0\sm k$, then there is an integer $m$ with $(\psi'-m[\tau])(K_0^\times)\subset pW$.
\end{minipage}
\end{equation}
For primitive $x$ the element $\dlog_A\psi'(x)$ lies in $\omega_A(Q_x)=F\dlog\tau(x)$, say $\dlog_A\psi'(x)=\theta'(x)\dlog\tau(x)$. As in the proof of Lemma~\ref{lem:theta}, $\theta'$ is constant on primitive elements (use $\psi'(xy)=\psi'(x)+\psi'(y)$ for transverse $x,y$),
lies in $l$, and lies in $\FF=\FF_p$ by Lemma~\ref{lem:R}; here one uses that a representative $(N,a')$ of $\psi'(x)\in W\sm0$ has $a'\in Q_x\cap F$, so that $a'$ is separable algebraic over $l(\tau(x))$ by Lemma~\ref{lem:pointsmodel}. Let $m\in\ZZ$ be a lift of $\theta'$. Then $\dlog_A(\psi'(x)-m[\tau x])=0$ for primitive $x$, so $\psi'(x)-m[\tau x]\in pW$ by Lemma~\ref{lem:W}(c).
For $x=\Frob^j(x_0)$ as in \eqref{eq:roots} the same element is $p^j$ times the one for $x_0$, and for $x\in k^\times$ it is zero. This proves \eqref{eq:lift}.

Starting with $\psi_0=\psi|_{K_0^\times}$ and applying \eqref{eq:lift} repeatedly we obtain integers $m_0,m_1,\dots$ and homomorphisms $\psi_j\colon K_0^\times\to W$ with $p\,\psi_{j+1}=\psi_j-m_j[\tau]$. Hence, with $r_j=\sum_{i<j}m_ip^i$,
\begin{equation}\label{eq:padic}
\psi=r_j[\tau]+p^j\psi_j\qquad\text{on }K_0^\times,\ \text{for all }j\ge0 .
\end{equation}
We now show that the $r_j$ converge $p$-adically to a rational number. Fix a primitive $x$, put $s=\tau(x)$, and let $(N,a)$ be a representative of $\psi(x)$. Then $a$ and $s$ are model generators of $Q_x$, so $a$ is separable algebraic over $l(s)$.
Let $E=l(s,a)$, $M=[E:l(s)]$ and $f=N_{E|l(s)}(a)\in l(s)^\times$. Since $ds\ne0$ in $\Omega$, the map $\Omega_{E|l}\to\Omega$ is injective, and therefore
\begin{equation}\label{eq:Epower}
E\cap L_A^{\,p^j}=E^{p^j}\qquad(j\ge0):
\end{equation}
if $y\in E$ is a $p$-th power in $L_A$ then $dy=0$, so $y\in E^p$ by \ref{S3}, and one proceeds by induction.
By \eqref{eq:padic}, $N p^j\psi_j(x)=[a\,s^{-Nr_j}]$. If $\psi_j(x)\neq0$ it has a representative $(N_j,g)$ with $g\in L_A$, and then $(as^{-Nr_j})^{N_j}=c\,g^{Np^j}$ with $c\in l^\times$; if $\psi_j(x)=0$ then $as^{-Nr_j}\in l$.
In both cases there is an integer $N_j$ prime to $p$ with $(as^{-Nr_j})^{N_j}\in l^\times\cdot E^{p^j}$, by \eqref{eq:Epower}. Taking norms,
\[ \bigl(f\,s^{-MNr_j}\bigr)^{N_j}\in l^\times\cdot l(s)^{p^j}\qquad\text{for all }j\geq0. \]
Write $f=c_0\,s^{\varepsilon}\prod_PP^{\varepsilon_P}$ with monic irreducible polynomials $P\neq s$. Unique factorization in $l[s]$ gives $\varepsilon_P\equiv0$ and $\varepsilon\equiv MNr_j$ modulo $p^j$, for all $j$.
Write $MN=p^\beta n'$ with $p\nmid n'$. Taking $j>\beta$ we see that $p^\beta$ divides $\varepsilon$, say $\varepsilon=p^\beta m'$, and
\[ n'r_j\equiv m'\pmod{p^{\,j-\beta}}\qquad\text{for all }j>\beta. \]
By \eqref{eq:padic}, $n'\psi-m'[\tau]=(n'r_j-m')[\tau]+n'p^j\psi_j$ maps $K_0^\times$ into $p^{\,j-\beta}W$ for all $j>\beta$, hence into $\bigcap_jp^jW=0$. So $\psi=\frac{m'}{n'}[\tau]$ on $K_0^\times$.

In both cases $r\ne0$ because $\psi\ne0$.
\end{proof}

\subsection{End of the proof}
By \eqref{eq:roots} and since $L$ is perfect, $\tau$ extends uniquely to a field embedding $\tau\colon K\to L$ commuting with Frobenius, and $\psi=r[\tau]$ holds on $K^\times$ because $V$ is uniquely divisible.
By Lemma~\ref{lem:tauinQ}, $\tau^{-1}(l)=k$. If $x\in K$ is algebraic over $k$, then $\tau(x)$ is algebraic over $l$, so $\tau(x)\in l$ and $x\in k$; thus $k$ is relatively algebraically closed.
It remains to prove uniqueness.

\begin{lemma}\label{lem:unique}
Let $\tau,\tau'\colon K\to L$ be field embeddings with $\tau^{-1}(l)=\tau'^{-1}(l)=k\neq K$, and $r,r'\in\QQ^\times$ with $r[\tau(x)]=r'[\tau'(x)]$ in $V$ for all $x\in K^\times$.
Then $\tau'=\Frob^j\circ\tau$ and $r'=p^{-j}r$ for some $j\in\Gamma$.
\end{lemma}

\begin{proof}
Write $r/r'=p^{j}m/m'$ with $j\in\Gamma$ and integers $m,m'$ prime to $p$ (if $p=1$, just $r/r'=m/m'$). Replacing $\tau$ by $\Frob^j\circ\tau$ and $r$ by $p^{-j}r$ we may assume $j=0$.
Let $x\in K\sm k$ and put $S=\tau(x)$, $S'=\tau'(x)$. The hypothesis for $x$ and for $1-x$ gives
\[ S^m=c\,S'^{m'},\qquad (1-S)^m=c'\,(1-S')^{m'}\qquad(c,c'\in l^\times). \]
As $m,m'$ are prime to $p$, $S$ and $S'$ are separable algebraic over each other; in $E=l(S,S')$ we have $\Omega_{E|l}=E\,dS$ with $dS'\ne0$, and taking logarithmic derivatives
\[ m\frac{dS}{S}=m'\frac{dS'}{S'},\qquad m\frac{dS}{1-S}=m'\frac{dS'}{1-S'} . \]
Dividing, $S/(1-S)=S'/(1-S')$, so $S=S'$. Hence $\tau'=\tau$ on $K\sm k$, therefore on $K$, and then $r=r'$.
\end{proof}

This completes the proof of Theorem~\ref{thm:A}. \qed


\section{Milnor \texorpdfstring{$K$}{K}-theory}\label{sec:K}

In this section $K$ and $L$ are perfect fields. We write $L^\times_\QQ=L^\times\otimes\QQ=\Kr_1(L)$ and denote the product in $\Kr_*$ by $\{\cdot,\cdot\}$.
Since $\Kr_*(L)$ is the quotient of the tensor algebra of $L^\times_\QQ$ by the ideal generated by the elements $[a]\otimes[1-a]$, a $\QQ$-linear map $B$ from $L^\times_\QQ\otimes L^\times_\QQ$ to a vector space factors through $\Kr_2(L)$ if and only if
$B([a],[1-a])=0$ for all $a\in L\sm\{0,1\}$.

\subsection{A tame symbol}

\begin{lemma}\label{lem:tame}
Let $v$ be a valuation of $L$ with valuation ring $O$, residue field $\kappa$ and value group $\Gamma_v$, and put $\Gamma_{v,\QQ}=\Gamma_v\otimes\QQ$, $\kappa^\times_\QQ=\kappa^\times\otimes\QQ$.
There is a $\QQ$-linear map
\[ T_v\colon\Kr_2(L)\to\Gamma_{v,\QQ}\otimes_\QQ\kappa^\times_\QQ\qquad\text{with}\qquad T_v\{c,b\}=-v(b)\otimes[\bar c]\quad\text{for all }c\in O^\times,\ b\in L^\times . \]
\end{lemma}

\begin{proof}
The valuation induces a surjection $v\colon L^\times_\QQ\to\Gamma_{v,\QQ}$ whose kernel $U$ is the image of $O^\times\otimes\QQ$, and the residue map induces $\mathrm{res}\colon U\to\kappa^\times_\QQ$.
Choose a $\QQ$-linear retraction $\pi\colon L^\times_\QQ\to U$ of the inclusion and put $\mathfrak u=\mathrm{res}\circ\pi$. Define
\[ B(\alpha,\beta)=v(\alpha)\otimes\mathfrak u(\beta)-v(\beta)\otimes\mathfrak u(\alpha). \]
This is bilinear and alternating, and has the stated value on $(c,b)$ with $c\in O^\times$, since then $v(c)=0$ and $\mathfrak u([c])=[\bar c]$. It remains to check $B([a],[1-a])=0$ for $a\in L\sm\{0,1\}$.
If $v(a)>0$, then $1-a$ is a unit with residue $1$, so $v(1-a)=0$ and $\mathfrak u([1-a])=0$.
If $v(a)<0$, then $[1-a]=[a]+[1-a^{-1}]$ in $L^\times_\QQ$ and $1-a^{-1}$ is a unit with residue $1$; hence $v(1-a)=v(a)$, $\mathfrak u([1-a])=\mathfrak u([a])$ and $B=0$.
If $v(a)=0$ and $v(1-a)=0$, both terms vanish. If $v(a)=0$ and $v(1-a)>0$, the first case applied to $1-a$ gives $B([1-a],[a])=0$.
\end{proof}

We also use the following standard fact: if $M\subset L$ is a subfield and $b\in L$ is transcendental over $M$, there is a valuation $v$ of $L$ which is trivial on $M$ and satisfies $v(b)>0$.
(Extend the $b$-adic valuation of $M(b)$ to $L$.) Its residue field contains $M$.

\begin{lemma}\label{lem:symbols}
Let $l\subset L$ be relatively algebraically closed.
\begin{enumerate}[label=\textup{(\alph*)}]
\item If $a,b\in L$ are algebraically independent over $l$, then $\{a,b\}\neq0$ in $\Kr_2(L|l)$.
\item If $c\in L^\times$ is algebraic over the prime field and not a root of unity, and $b\in L$ is transcendental over the prime field, then $\{c,b\}\ne0$ in $\Kr_2(L)$.
\end{enumerate}
\end{lemma}

\begin{proof}
(a) Let $M$ be the relative algebraic closure of $l(a)$ in $L$, and $v$ a valuation trivial on $M$ with $v(b)>0$. Composing $T_v$ with the projection $\kappa^\times_\QQ\to(\kappa^\times/l^\times)_\QQ$ gives a map $T'_v$ which kills $\{c,y\}$ for $c\in l^\times$,
hence factors through $\Kr_2(L|l)$. Now $T'_v\{a,b\}=-v(b)\otimes[\bar a]$, and $[\bar a]\neq0$ in $(\kappa^\times/l^\times)_\QQ$ because no power of $a$ lies in $l$.
(b) Let $l_0$ be the algebraic closure of the prime field in $L$ and $v$ a valuation trivial on $l_0$ with $v(b)>0$. Then $T_v\{c,b\}=-v(b)\otimes[\bar c]\ne0$.
\end{proof}

In terms of $V_{L|l}=\Kr_1(L|l)$, part (a) says: if $v_1,v_2\in V_{L|l}$ satisfy $\{v_1,v_2\}=0$ in $\Kr_2(L|l)$, then $v_1,v_2$ are algebraically dependent or one of them is $0$.

\subsection{Proof of Theorem~\ref{thm:B}}
Let $\psi=\varphi_1\circ[\,\cdot\,]\colon K^\times\to V_{L|l}$. Since $\varphi_1$ is injective and $k$ is relatively algebraically closed, $\ker\psi=k^\times$ by \ref{V1}.
For $x\in K\sm k$ we have $\{\psi(x),\psi(1-x)\}=\varphi_2\{x,1-x\}=0$, and both elements are non-zero, so they are algebraically dependent by Lemma~\ref{lem:symbols}(a). Thus \ref{A1}--\ref{A4} hold, and Theorem~\ref{thm:A} gives
$\tau$ and $r$ with $\varphi_1=r\,\tau_*$ in degree $1$. As both $\varphi$ and $r^{\deg}\circ\tau_*$ are homomorphisms of algebras generated in degree $1$, they agree. Uniqueness is part of Theorem~\ref{thm:A}. \qed

\subsection{Proof of Theorem~\ref{thm:C}}
Let $\psi$ be the composition of $\varphi_1\circ[\,\cdot\,]\colon K^\times\to L^\times_\QQ$ with the projection to $V=V_{L|l_0}$, and let $C_L\subset L^\times_\QQ$ be the kernel of this projection, the image of $l_0^\times\otimes\QQ$.
Put $k'=\{0\}\cup\ker\psi$. As in the proof of Theorem~\ref{thm:B}, Lemma~\ref{lem:symbols}(a) gives:
\begin{equation}\label{eq:dep}
\begin{minipage}{0.86\textwidth}
for every $x\in K\sm\{0,1\}$, the elements $\psi(x)$ and $\psi(1-x)$ of $V$ are algebraically dependent, or one of them is $0$.
\end{minipage}
\end{equation}

\begin{lemma}\label{lem:kprime}
$k'$ is a subfield of $K$.
\end{lemma}

\begin{proof}
$\ker\psi$ is a subgroup containing $-1$, so it suffices to show: if $z\in\ker\psi$ and $z\ne1$, then $1-z\in\ker\psi$.

\emph{(i) $z$ is not a root of unity.} Then $\varphi_1[z]$ is a non-zero element of $C_L$, say $\varphi_1[z]=\frac1N[c]$ with $c\in l_0$ not a root of unity. Write $\varphi_1[1-z]=\frac1{N'}[b]$. Then $\{c,b\}=NN'\varphi_2\{z,1-z\}=0$,
so $b$ is algebraic over the prime field by Lemma~\ref{lem:symbols}(b), i.e.\ $1-z\in\ker\psi$.

\emph{(ii) $z=\zeta$ is a root of unity and $\chr K>0$.} Then $1-\zeta$ is algebraic over a finite field, hence a root of unity, and $\varphi_1[1-\zeta]=0$.

\emph{(iii) $z=\zeta$ is a root of unity and $\chr K=0$.} Then $\chr L=0$ by \ref{A4}. Let $\Delta=\dlog\circ\psi\colon K^\times\to\Omega_{L|l_0}$, where $\dlog\colon V\to\Omega_{L|l_0}$ is the injective map of Lemma~\ref{lem:W}(c) (with $L_A=L$); so $\ker\Delta=\ker\psi$. By \ref{S2}, algebraically dependent elements of $V$ have proportional images under $\dlog$, and algebraically independent ones have linearly independent images. Suppose $c=1-\zeta\notin\ker\psi$, so $\Delta(c)\ne0$.
We claim that there is $x\in K$, transcendental over $\QQ$, with $\Delta(x)\notin L\,\Delta(c)$. By \ref{A3} there is $x_1$ with $\Delta(x_1)\notin L\Delta(c)$. If $x_1$ is transcendental we are done.
Otherwise let $z_0\in K\sm k_0$; then $z_0$ and $z_0x_1$ are transcendental, and one of $\Delta(z_0)$, $\Delta(z_0x_1)=\Delta(z_0)+\Delta(x_1)$ is not in $L\Delta(c)$.

Fix such an $x$; then $x$ and $cx$ are not in $\ker\psi$. By \eqref{eq:dep}, $\Delta(1-x)=\mu\Delta(x)$ and $\Delta(1-cx)=\mu'\Delta(cx)$ with $\mu,\mu'\in L$. Let $u=c(1-x)/(c-1)$. Since $c-1=-\zeta$ and $1-u=(1-cx)/\zeta$,
\[ \Delta(u)=\Delta(c)+\mu\,\Delta(x)\neq0,\qquad \Delta(1-u)=\Delta(1-cx)=\mu'\bigl(\Delta(c)+\Delta(x)\bigr), \]
and these are proportional by \eqref{eq:dep}. Hence $\mu'=0$ or $\mu=1$.

If $\mu'=0$, then $1-cx\in\ker\psi$. It cannot be a non-root of unity, since then (i) would give $cx=1-(1-cx)\in\ker\psi$. So $1-cx$ is a root of unity and $x$ is algebraic over $\QQ$.
If $\mu=1$, then $\Delta((1-x)/x)=0$, so $1-1/x\in\ker\psi$; by (i) it is a root of unity, as otherwise $1/x\in\ker\psi$. Again $x$ is algebraic over $\QQ$.
Both cases contradict the choice of $x$. Hence $1-\zeta\in\ker\psi$.
\end{proof}

By Lemma~\ref{lem:kprime} and \eqref{eq:dep}, the homomorphism $\psi\colon K^\times\to V_{L|l_0}$ satisfies \ref{A1} and \ref{A2} with $k=k'$; \ref{A3} and \ref{A4} hold by assumption.
Theorem~\ref{thm:A} gives a field embedding $\tau\colon K\to L$ with $\tau^{-1}(l_0)=k'$ and $r\in\QQ^\times$ with $\psi=r[\tau]$ in $V_{L|l_0}$.
An element of $K$ is algebraic over the prime field if and only if its image under $\tau$ is, so $k'=\tau^{-1}(l_0)=k_0$.

It remains to show that $\varphi_1[x]=r[\tau(x)]$ holds in $L^\times_\QQ$, not only modulo $C_L$. Let $\chi(x)=\varphi_1[x]-r[\tau(x)]\in C_L$; this is a homomorphism $K^\times\to C_L$.
If $\chr L>0$ then $C_L=0$. Let $\chr L=0$ and $x\in K\sm k_0$, so $s=\tau(x)$ is transcendental over $l_0$. Choose a valuation $v$ of $L$, trivial on $l_0$, with $v(s)>0$. Then $v(1-s)=0$, and $v$ vanishes on $C_L$.
Expanding $0=T_v\,\varphi_2\{x,1-x\}=T_v\{r[s]+\chi(x),\,r[1-s]+\chi(1-x)\}$ with the formula of Lemma~\ref{lem:tame}, all terms vanish except
\[ T_v\{r[s],\chi(1-x)\}=r\,v(s)\otimes\overline{\chi(1-x)}, \]
where the bar denotes the injection $C_L\to\kappa^\times_\QQ$. Hence $\chi(1-x)=0$ for all $x\in K\sm k_0$, i.e.\ $\chi$ vanishes on $K\sm k_0$, and then also on $k_0^\times$.
So $\varphi_1=r\,\tau_*$, and $\varphi=r^{\deg}\circ\tau_*$ as in the proof of Theorem~\ref{thm:B}. Uniqueness follows from Theorem~\ref{thm:A}. \qed

\subsection{Proof of Theorem~\ref{thm:D}}
We may replace $K,L$ by their perfect closures. Let $\varphi$ be an isomorphism. Both $\varphi_1$ and $\varphi_1^{-1}$ are surjective, so \ref{A3} holds for $\varphi$ and for $\varphi^{-1}$ with $d=3$, and with $d=2$ for the one whose target has characteristic $0$.
If $\chr K=0<\chr L$, Theorem~\ref{thm:C} applies to $\varphi^{-1}$ and gives an embedding $L\to K$, which is absurd; similarly $\chr L=0<\chr K$ is impossible. In particular, if one of the fields has characteristic $0$ so does the other.
So \ref{A4} holds in both directions, and Theorem~\ref{thm:C} gives embeddings $\tau\colon K\to L$, $\tau'\colon L\to K$ and $r,r'\in\QQ^\times$ with $\varphi=r^{\deg}\tau_*$ and $\varphi^{-1}=r'^{\deg}\tau'_*$.
Then $rr'\,[\tau'\tau(x)]=[x]$ in $K^\times_\QQ$ for all $x$, and the uniqueness in Theorem~\ref{thm:A}, applied to the canonical map $K^\times\to V_{K|k_0}$, shows that $\tau'\tau$ is a power of Frobenius. The same holds for $\tau\tau'$.
Hence $\tau$ is an isomorphism. \qed

\section{Remarks}\label{sec:remarks}

\begin{remark}[The hypotheses]
\ref{A1} cannot be dropped: if $v$ is a valuation of $K$ and $\chi\colon\Gamma_v\to V_{L|l}$ is any homomorphism, then $\psi=\chi\circ v$ satisfies $\psi(1-x)\in\{0,\psi(x)\}$ whenever $\psi(x)\ne0$, so the dependence condition in \ref{A2} holds trivially.
More generally one expects that without injectivity $\psi$ is governed by a place of $K$ rather than an embedding; this is not treated here.
\ref{A3} excludes the homomorphisms whose image lies in a single point. Hypothesis \ref{A4} is used only to know that $K^\times$ has no non-zero homomorphism to $\ZZ$ when $\chr L>0$.
\end{remark}

\begin{remark}[Open points]
(1) In positive characteristic the proof needs three $\psi$-independent elements, for the seven-point configuration of Lemma~\ref{lem:cocycle}. I do not know whether Theorem~\ref{thm:A} holds with $d=2$ in positive characteristic.
(2) I do not know whether there is a homomorphism $\psi$ satisfying \ref{A1}--\ref{A3} with $\chr K=0<\chr L$. Hypothesis \ref{A4} enters only through \eqref{eq:nohom}.
(3) Theorem~\ref{thm:C} assumes $K\neq k_0$. For fields algebraic over the prime field, $\Kr_*$ vanishes in positive degrees in positive characteristic, and in characteristic $0$ the question is of a different nature.
\end{remark}

\begin{remark}[On formalization]
The logical structure is as follows. Section~\ref{sec:conv} is standard field theory. Sections~\ref{sec:levels}--\ref{sec:quadrangle} concern a single extension $L|l$ in characteristic $p$ and are used only through
Lemma~\ref{lem:normalized} (existence and uniqueness of level triples), Lemmas~\ref{lem:pointsmodel}, \ref{lem:lax} (points in a model) and Lemma~\ref{lem:quadrangle}; they enter the proof of Theorem~\ref{thm:A} only in Lemmas~\ref{lem:gauge}, \ref{lem:cocycle} and \ref{lem:theta}.
In characteristic $0$ the proof of Theorem~\ref{thm:A} consists of Lemmas~\ref{lem:onevar}--\ref{lem:tconst}, Proposition~\ref{prop:tau}, Lemma~\ref{lem:R}, Lemmas~\ref{lem:tauinQ}, \ref{lem:theta}, the first case of Proposition~\ref{prop:rational} and Lemma~\ref{lem:unique}, with $F=L$, $K_0=K$, $\mathcal N=k^\times$, $\Delta=\dlog\circ\psi$;
the only facts about differentials needed are that algebraically dependent elements have proportional differentials and that $da=0$ implies $a\in l$.
Section~\ref{sec:K} needs, in addition, the extension of valuations.
\end{remark}

\begin{thebibliography}{9}
\bibitem{BT} F.~Bogomolov and Yu.~Tschinkel, \emph{Milnor $K_2$ and field homomorphisms}, Surveys in Differential Geometry \textbf{13} (2009), 223--244.
\bibitem{T} A.~Topaz, \emph{Algebraic dependence and Milnor $K$-theory}, preprint, arXiv:2211.14665.
\end{thebibliography}

\end{document}
